\indent ಮೇಲೆ ಕೊಟ್ಟಿರುವ ಉದಾಹರಣೆಗಳಲ್ಲಿ ಓಡು, ಹಾಡು, ತಿನ್ನು, ಆಡು ಕ್ರಿಯಾಪ್ರಕೃತಿಗಳು ಅಥವಾ ಧಾತುಗಳು ; “ಉತ್ತ” ಎಂಬುದು ಕಾಲಸೂಚಕ ಪ್ರತ್ಯಯ ಮತ್ತು ಇದೆ, ಆಳೆ, ಅದೆ, ಆನೆ ಆಖ್ಯಾತ ಪ್ರತ್ಯಯಗಳು. ಸೂತ್ರ ರೂಪದಲ್ಲಿ ಹೇಳುವುದಾದರೆ,\par

\indent ಕ್ರಿಯಾ = ಕ್ರಿಯಾ + ಕಾಲಸೂಚಕ + ಆಖ್ಯಾತ\par

\indent ಪದ + ಪ್ರಕೃತಿ + ಪ್ರತ್ಯಯ + ಪ್ರತ್ಯಯ ವಿಶೇಷ ಸೂಚನೆ : ಕೆಲವು ನಾಮಪ್ರಕೃತಿಗಳನ್ನು ಮತ್ತು ಅನುಕರಣ ಶಬ್ದಗಳನ್ನು ಧಾತುಗಳಂತೆ ಬಳಸುವುದುಂಟು. ನಾಮಪ್ರಕೃತಿಯ ಮೇಲಾಗಲೀ ಅನುಕರಣ ಶಬ್ದದ ಮೇಲಾಗಲೀ “ಇಸು” ಪ್ರತ್ಯಯ ಸೇರಲು ಅದು ಧಾತು ವೆನಿಸುವುದು. ಅಂತಹ ಧಾತುಗಳಿಗೆ ಪ್ರತ್ಯಯಾಂತ ಧಾತುಗಳೆಂದು ಹೆಸರು.\par

\vspace{3pt plus 2pt}\noindent\hspace*{1.2em}\kbvlabel{ಉದಾ}\par\vspace{1pt}

\begin{kbvderiv}
\vspace{4pt plus 3pt minus 1pt}
\begin{center}\begin{adjustbox}{max width=\textwidth}
{\renewcommand{\arraystretch}{1.22}\setlength{\tabcolsep}{4pt}
\begin{tabular}{l c l}
ಕನ್ನಡ + ಇಸು & {\color{kbvgrey}$=$} & \kbvresult{ಕನ್ನಡಿಸು}{\color{kbvgrey} (ಪ್ರತ್ಯಯಾಂತ ಧಾತು)} \\
ಉತ್ತರ + ಇಸು & {\color{kbvgrey}$=$} & \kbvresult{ಉತ್ತರಿಸು}{\color{kbvgrey} ( ೨ )} \\
“ಪ್ರಯತ್ನ + ಇಸು & {\color{kbvgrey}$=$} & \kbvresult{ಪ್ರಯತ್ನಿಸು}{\color{kbvgrey} ( ೨ )} \\
ಸಿದ್ಧಿ + ಇಸು & {\color{kbvgrey}$=$} & \kbvresult{ಸಿದ್ಧಿಸು}{\color{kbvgrey} ( ೨ )} \\
ಥಳಥಳ + ಇಸು & {\color{kbvgrey}$=$} & \kbvresult{ಥಳಥಳಿಸು}{\color{kbvgrey} ( ೨ )} \\
ಘಮಘಮ + ಇಸು & {\color{kbvgrey}$=$} & \kbvresult{ಘಮಘಮಿಸು}{\color{kbvgrey} ( ೨ )} \\
\end{tabular}}
\end{adjustbox}\end{center}
\vspace{4pt plus 3pt minus 1pt}
\end{kbvderiv}

\indent ಪ್ರತ್ಯಯಾಂತ ಧಾತುಗಳು ಆಖ್ಯಾತ ಪ್ರತ್ಯಯಗಳನ್ನು ಹೊಂದಿ ಕ್ರಿಯಾಪದಗಳಾಗುವುವು.\par

\begin{kbvderiv}
\vspace{4pt plus 3pt minus 1pt}
\begin{center}\begin{adjustbox}{max width=\textwidth}
{\renewcommand{\arraystretch}{1.22}\setlength{\tabcolsep}{4pt}
\begin{tabular}{l c l}
ಉದಾ : ಕನ್ನಡಿಸು + ಉತ್ತ + ಆನೆ & {\color{kbvgrey}$=$} & \kbvresult{ಕನ್ನಡಿಸುತ್ತಾರೆ} \\
ಸಿದ್ಧಿಸು + ಇದೆ & {\color{kbvgrey}$=$} & \kbvresult{ಸಿದ್ಧಿಸಿತು} \\
ಥಳಥಳಿಸು + ಉದು & {\color{kbvgrey}$=$} & \kbvresult{ಥಳಥಳಿಸುವುದು} \\
\end{tabular}}
\end{adjustbox}\end{center}
\vspace{4pt plus 3pt minus 1pt}
\end{kbvderiv}

\kbvsechead{ಅಕರ್ಮಕ ಸಕರ್ಮಕ ಧಾತುಗಳು :}

\begin{kbvrule}
\noindent\kbvbullet\hspace{0.5em}ಧಾತುಗಳಲ್ಲಿ ಎರಡು ವಿಧ. ಆ ಎರಡು ವಿಧಗಳು ಯಾವುವೆಂದರೆ,
\end{kbvrule}

\vspace{2.5pt plus 3pt minus 0.5pt}\noindent\hspace*{1.2em}\kbvnum{(೧)}\hspace{0.55em}ಅಕರ್ಮಕ ಧಾತು ಮತ್ತು (೨) ಸಕರ್ಮಕ ಧಾತು.\par
